% !TEX program = lualatex
\documentclass[11pt,a4paper]{article}
\usepackage{icat-common}
\usepackage{hyperref}

%-------------------------------------------------
%  独自マクロ定義
%-------------------------------------------------
\newcommand{\lhagma}{\mathcal{L}}
\newcommand{\sem}{\mathsf{Sem}}
\newcommand{\refe}{\mathsf{Ref}}
\newcommand{\sense}{\mathsf{Sense}}
\newcommand{\ctx}{\Gamma}
\newcommand{\botext}{\bot_{\mathrm{ext}}}
\newcommand{\botint}{\bot_{\mathrm{int}}}
\newcommand{\lext}{\lhagma_{\mathrm{ext}}}
\newcommand{\lint}{\lhagma_{\mathrm{int}}}
\newcommand{\yon}{\mathsf{Y}}

%-------------------------------------------------
%  定理環境の設定
%-------------------------------------------------
\theoremstyle{definition}
\newtheorem{definition}{定義}[section]
\newtheorem{thesis}{命題}[section]
\newtheorem{theorem}{定理}[section]
\newtheorem{remark}{補足}[section]

\begin{document}

\title{残余意味論：意味と指示の証明論的再構成}
\author{千葉将臣}
\date{2026-06-26}
\maketitle

\begin{abstract}
本稿は、意味（sense）と指示（reference）の関係を、対象の直接的同一性ではなく、外部的指定と内部的特徴付けの対応として再構成する試みである。通常、指示は対象が何を指すかを与え、意味はその対象がどのような仕方で与えられるかを与えるとされる。本稿では、この二分法に対して、対象水準の記述を排除した後にもなお残る「残余」$\lhagma$ を導入し、指示を $\lext$、意味を $\lint$ として形式化する。そのうえで、直観主義論理における矛盾元 $\bot$ が、外部から導入される定数でありながら、局所的な矛盾文脈から体系内で導出される点に注目する。この二重性を模型として、残余 $\lhagma$ を、外部的に指定される指示と内部的に構成される意味が対応する特異点として解釈する。最後に、この対応を米田の補題そのものではなく、対象がその文脈的振る舞いによって特徴付けられるという米田型の類比として位置付ける。
\end{abstract}

\section{序論}

意味と指示の区別は、言語哲学および論理意味論における基本問題である。指示は表現が何を指すかに関わり、意味はその指示対象がどのような認識的・論理的経路を通じて与えられるかに関わる。本稿の目的は、この区別を証明論的に再構成することである。

本稿で中心となる概念は「残余」$\lhagma$ である。残余とは、対象水準の述語、構成、または指示可能な内容を排除した後にも、なお体系上の役割として残るものを指す。残余は、単なる対象ではなく、外部的に指定される側面と、内部的な推論過程を通じて特徴付けられる側面をあわせ持つ。本稿では前者を $\lext$、後者を $\lint$ と表記する。

本稿の主張は、意味と指示の関係を、単純な同一性としてではなく、次の対応として捉える点にある。
\[
\lext \cong \lint .
\]
この式は、意味と指示が区別なく同じであることを意味しない。むしろ、外部的に指定された残余の指示が、体系内部の推論過程によって意味として回収されることを表す。

\section{意味・指示・残余}

\begin{definition}[指示]
指示 $\refe(A)$ とは、表現 $A$ が外部的に何を指すかを表す側面である。本稿では、残余 $\lhagma$ の指示的側面を $\lext$ と書く。
\end{definition}

\begin{definition}[意味]
意味 $\sense(A)$ とは、表現 $A$ がどのような推論過程、排他過程、または文脈的振る舞いを通じて特徴付けられるかを表す側面である。本稿では、残余 $\lhagma$ の意味的側面を $\lint$ と書く。
\end{definition}

\begin{definition}[残余]
残余 $\lhagma$ とは、対象水準の述語や構成を排除した後にも、体系内でなお区別可能な役割として残るものを指す。残余は、対象的内容を持つ通常の指示対象ではなく、外部的指定と内部的特徴付けの対応点として定義される。
\end{definition}

以上の定義により、残余は次の二側面を持つ。
\[
\lhagma = (\lext, \lint).
\]
ここで $\lext$ は「何が残るか」を表し、$\lint$ は「それがどのように残るか」を表す。したがって、残余意味論における中心問題は、$\lext$ と $\lint$ がどのような条件のもとで対応するかである。

\section{直観主義論理における矛盾元の二重性}

直観主義論理において、$\bot$ は通常、矛盾または構成不可能性を表す定数として導入される。この意味では、$\bot$ は外部から体系の語彙に与えられる基礎記号である。これを $\botext$ と書く。

一方で、$\bot$ は局所的な矛盾文脈から体系内で導出される。たとえば、文脈 $\ctx$ が命題 $P$ とその否定 $\neg P$、すなわち $P \to \bot$ を含むとき、含意除去により次が得られる。
\[
\frac{P \qquad P \to \bot}{\bot}
\]
この場合の $\bot$ は、外部から与えられた記号ではなく、体系内部の推論過程によって到達される結論である。これを $\botint$ と書く。

したがって、直観主義論理における $\bot$ には、少なくとも次の二つの側面がある。
\[
\botext \qquad \text{および} \qquad \botint .
\]
本稿では、この二つが同一の記号 $\bot$ において一致することを、次の作業上の同型関係として表す。
\[
\botext \cong \botint .
\]
この同型は、無条件に $\vdash \bot$ が成り立つことを意味しない。むしろ、$\bot$ が外部から導入される定数であると同時に、矛盾的な局所文脈のもとで体系内から到達される特異点であることを表す。

\section{残余意味論の基本命題}

前節の $\bot$ の二重性は、残余 $\lhagma$ の構造を理解するための模型を与える。すなわち、残余もまた、外部的に指定される側面と、内部的に導出または特徴付けられる側面を持つ。

\begin{definition}[外部的残余]
外部的残余 $\lext$ とは、対象水準の記述を排除した後に、なお「残るもの」として指定される指示的側面である。
\end{definition}

\begin{definition}[内部的残余]
内部的残余 $\lint$ とは、対象水準の述語を排除する過程、矛盾を局所化する過程、または推論を正規化する過程を通じて、体系内部で特徴付けられる意味的側面である。
\end{definition}

\begin{thesis}[残余意味論の基本命題]
残余 $\lhagma$ は、外部的に指定される指示 $\lext$ と、内部的に特徴付けられる意味 $\lint$ が対応する点として定式化される。
\[
\lext \cong \lint .
\]
\end{thesis}

この命題は、意味と指示の単純な同一視ではない。指示は外部的指定として与えられ、意味は内部的構成として与えられる。両者は異なる記述水準に属するが、残余 $\lhagma$ において対応する。この対応が成立する点で、残余は通常の対象ではなく、意味と指示の接続を担う形式的特異点となる。

\section{米田型対応としての残余}

圏論における米田の補題は、対象がその対象へ向かう射、またはその対象から出る射の全体によって特徴付けられることを示す。本稿は米田の補題そのものを直接適用するものではない。しかし、残余意味論における $\lext \cong \lint$ は、外部的に指定された対象が、その対象をめぐる文脈的振る舞いによって内部的に特徴付けられるという点で、米田型の類比を持つ。

この類比を明示するため、文脈 $\ctx$ から残余 $\lhagma$ への到達可能性を、形式的に次のように表す。
\[
\yon(\lhagma)(\ctx) := \mathrm{Hom}(\ctx, \lhagma).
\]
ここで $\mathrm{Hom}(\ctx, \lhagma)$ は、文脈 $\ctx$ から残余 $\lhagma$ へ至る推論、排他、または正規化の経路を表す記号である。このとき、内部的残余 $\lint$ は、単独の対象としてではなく、文脈から残余へ至る振る舞いの総体として理解される。
\[
\lint \simeq \yon(\lhagma).
\]
したがって、外部的に指定される $\lext$ と、文脈的振る舞いとして特徴付けられる $\lint$ が対応することは、残余が外部指定と内部構成の交差点であることを意味する。

\begin{remark}[米田型という限定]
本稿でいう米田型対応は、圏論上の米田の補題を厳密に適用する主張ではない。ここで意図されるのは、対象の同一性を、その対象をめぐる文脈的振る舞いによって特徴付けるという方法論的類比である。
\end{remark}

\section{意味と指示の再構成}

残余意味論の観点では、意味と指示の関係は次のように整理される。
\[
\refe(\lhagma) = \lext,
\qquad
\sense(\lhagma) = \lint .
\]
したがって、意味と指示の関係は、次の対応として表される。
\[
\refe(\lhagma) \cong \sense(\lhagma).
\]
ただし、この式は意味と指示が区別なく同じであることを主張しない。むしろ、残余においては、外部的指示が内部的意味構成によって回収されるという特別な構造が成立することを表す。

この再構成により、意味は単なる主観的把握様式ではなく、体系内部の推論過程として扱われる。また、指示は単なる外部対象ではなく、内部的意味構成によって回収可能な残余として扱われる。したがって、残余意味論は、意味と指示の関係を、対象と記述の関係ではなく、外部指定と内部導出の関係として定式化する。

\section{空論との関係}

本稿の議論は、空論に依存するものではない。しかし、空論における中道不動点 $\dfix$ と残余 $\lhagma$ の関係は、残余意味論の一つの応用例を与える。空論では、対象水準の差異を消去した後に残る真理保存則が、基底状態として定式化される。このとき、残余 $\lhagma$ は、外部的には「消去後に残るもの」として指定され、内部的には推論の帰着過程を通じて特徴付けられる。

したがって、空論における
\[
\lhagma \cong \dfix
\]
という対応は、残余意味論の観点からは、外部的指示と内部的意味構成の一致例として解釈できる。ただし、本稿の主張は空論の内部に限定されない。残余意味論は、意味と指示の関係一般を、証明論的・文脈論的に再記述するための独立した枠組みである。

\section{結論}

本稿では、意味と指示の関係を、残余 $\lhagma$ を介して証明論的に再構成した。まず、直観主義論理における $\bot$ が、外部から導入される定数でありながら、局所的矛盾文脈のもとでは体系内部から導出されることを確認した。この二重性を模型として、残余 $\lhagma$ を、外部的に指定される指示 $\lext$ と、内部的に特徴付けられる意味 $\lint$ が対応する点として定式化した。

この構成により、意味と指示の関係は、単純な同一性ではなく、外部指定と内部導出の対応として理解される。さらに、この対応は、米田の補題そのものではないが、対象をその文脈的振る舞いによって特徴付けるという意味で米田型の類比を持つ。以上により、残余意味論は、意味と指示の関係を、論理・証明論・文脈性の観点から再考するための枠組みを与える。

\begin{thebibliography}{9}
\bibitem{frege}
Gottlob Frege.
\textit{"Uber Sinn und Bedeutung}.
Zeitschrift f\"ur Philosophie und philosophische Kritik, 100, 1892.
本稿では、意味と指示の古典的区別を参照する基礎文献として用いる。

\bibitem{heyting}
A. Heyting.
\textit{Intuitionism: An Introduction}.
North-Holland, 1956.
本稿では、直観主義論理および矛盾元の形式化の参照文献として用いる。

\bibitem{yoneda}
Nobuo Yoneda.
On Ext and exact sequences.
\textit{Journal of the Faculty of Science, University of Tokyo, Section I}, 8, 1960.
本稿では、対象を文脈的振る舞いによって特徴付ける方法論的類比の参照として用いる。
\end{thebibliography}

\end{document}
